% chapters/01_introduction.tex - 第1章 前言

\chapter{前言}

\section{问题的提出}

抗阻训练是维持骨骼肌质量、神经肌肉机能与心代谢健康的核心运动方式，已被美国运动医学学会与世界卫生组织共同列为成年人健康促进的基础训练形式（\citep{ACSM2011}; \citep{WHO2020}）。近年来，随着全民健身战略与《健康中国2030》规划纲要的持续推进，大众群体对抗阻训练的科学认知与参与意愿显著提升，参与规模呈现进一步扩大态势。然而，参与规模的扩大并未伴随专业指导资源的同步供给。当前社会体育指导员队伍的专项技能多集中于传统健身项目（\citep{YangNan2025}），而商业化私人教练服务的费用门槛使多数普通训练者难以获得持续的专业指导。这一供需失衡导致大众抗阻训练者普遍处于缺乏专业监督的自我指导状态，需独立完成负荷设定、组数与次数安排以及组间恢复判断等核心训练决策。

自我训练场景中，训练者独立面对两类核心决策挑战：当日训练负荷的科学设定与训练过程中动作技术规范性的持续监控。就负荷设定而言，训练者主要依赖训练前的固定处方或即时主观感受进行判断，缺乏对当日神经肌肉状态的客观量化参照，因而难以准确识别负荷是否偏离目标刺激区间，也难以及时察觉疲劳累积带来的过度训练风险。就动作监控而言，训练者的本体感觉与视角局限使其难以在负荷增大与疲劳累积过程中及时察觉动作代偿与技术失范\citep{GwynneCurran2014}。这两类挑战在训练过程中存在交互效应：负荷设定失准会加剧动作技术失范的概率，而动作技术失范则会同时降低训练效益并增大骨骼肌肉系统的损伤风险。要在无专业监督条件下缓解这一双重挑战，训练者需要在训练过程中获取相对可靠的反馈信息以动态调节负荷与动作决策。反馈机制的科学性与可及性因此成为大众抗阻训练科学化的核心议题。

围绕反馈机制的构建，训练科学先后发展出两条方法学路径。第一条路径以主观感知反馈为基础，其代表性工具为主观用力程度量表\citep{Borg1982}与储备重复次数评估（\citep{Zourdos2016}），训练者依据即时用力感与对剩余能力的评估调节训练负荷。这一路径的方法学优势在于低成本与高可及性，其局限在于精度依赖训练经验，缺乏外部客观校准。第二条路径以客观测量反馈为基础，其代表性方法为基于速度的训练（velocity-based training, VBT），通过实时监测杠铃向心速度将训练强度锚定在生物力学参数之上\citep{GonzalezBadillo2010}。这一路径的方法学优势在于测量精度与决策客观性，其局限在于对专业硬件设备的高度依赖。传统VBT实现方案依赖线性位置传感器等专业设备，单套价格通常在数千美元级别，这一成本壁垒使其应用场景长期局限于竞技训练与科研实验室。两条路径在较长时期内相对独立地发展，各自积累证据基础，各自服务于不同的应用场景。

近年来，两条路径均出现了实质性变化。在主观反馈侧，基于RIR的自主调节在方法学证据上持续积累，包括RIR训练与力竭训练在肌肥大适应上的等效性验证（\citep{Refalo2024}）。基于这一证据基础，美国运动医学学会于2026年发布了距上一版本相隔17年的抗阻训练立场声明更新，正式将RIR取代传统的1RM百分比确立为强度处方的首选指标\citep{Currier2026}。这一更新意味着主观反馈路径在权威指南层面的方法学地位实质性提升。在客观反馈侧，深度学习姿态估计技术的成熟为VBT的低成本实现提供了新的技术路径，基于智能手机摄像头与轻量级姿态估计模型的移动端方案在原理上已具备完成杠铃速度测量任务的技术条件（\citep{Bazarevsky2020}）。这一技术演进为客观反馈路径向大众场景的下沉提供了成本与设备可及性上的可能。两条路径的近期变化在时间上高度重合：主观反馈路径在权威指南层面获得方法学地位的提升，客观反馈路径在技术层面获得成本可及性的改善。两条路径的这一同步演进使得一个具有现实意义的科学问题浮现：在规范化RIR训练已具备科学性、权威性与可及性的前提下，客观速度反馈是否仍具有值得投入额外成本的适应性增量？

围绕这一核心问题，本研究提出两层递进的追问。第一层追问关乎干预的必要性，即客观速度反馈相较于规范化RIR反馈能否在训练执行质量、疲劳控制以及体能适应上产生具有实践意义的增量差异。第二层追问关乎工具的可行性，即若客观反馈确有增量价值，基于计算机视觉的移动端方案能否在真实训练场景中提供接近专业硬件的测量能力与实用价值。为系统回应这两层追问，本研究构建了由三个阶段组成的递进式研究方案，研究一通过测量学效度验证，检验自研移动端AI系统与专业硬件在核心运动学参数测量上的一致性，为后续研究提供工具层面的方法学基础，研究二采用两组平行随机对照预试验设计，在真实训练场景下比较基于客观速度反馈与基于规范化RIR反馈的两种自主调节方案在4周下肢抗阻训练中的可行性、执行质量与初步适应差异。第三阶段采用半结构化深度访谈与主题分析，从主观体验层面揭示两种反馈路径下训练者的决策逻辑与工具接受机制。三个阶段依次回应工具效度、干预效果与用户体验三个层面的证据缺口，共同为大众抗阻训练场景下反馈方式的合理选择提供实证参照。

\section{选题的目的意义}

\subsection{选题目的}

本研究围绕1.1节提出的两层追问，以青年男性抗阻训练者为研究对象，通过工具效度验证、预试验性随机对照训练干预与半结构化深度访谈相结合的三阶段设计，系统评估基于客观速度反馈与基于规范化RIR反馈两条自主调节路径在真实抗阻训练场景中的相对表现。具体研究目的包括以下四个方面。第一，明确自研移动端AI系统在杠铃深蹲平均向心速度与反向运动跳高度测量中的一致性水平、系统偏差方向与绝对误差特征，回答基于计算机视觉的低成本方案在测量学层面能否达到支持真实训练监控应用的效度基础。第二，评估基于客观速度反馈的抗阻训练方案在真实场景中的招募保留率、训练完成率、速度目标命中率、自动调节规则执行稳定性与不良事件发生情况，对照预设推进标准判断该方案向后续大规模确证性随机对照试验推进的可行性水平。第三，比较两条反馈路径在训练负荷结构、内部负荷反应与初步体能适应上的差异特征，获取两组在深蹲最大力量、下肢爆发力与训练自我效能上的效应量点估计与置信区间，为后续正式随机对照试验的样本量计算提供先验参数。第四，从行为决策层面揭示训练者对客观速度反馈与主观RIR反馈的接受度差异，识别不同训练经验受试者在自主权让渡、算法信任构建与工具持续使用意愿上的分化特征，为数字化训练工具的用户导向优化提供质性证据。

\subsection{选题意义}

\subsubsection{理论意义}

本研究的理论贡献主要体现在两个维度。第一，本研究将抗阻训练自主调节领域的核心对比框架从"VBT vs. 传统百分比训练"推进至"VBT vs. 规范化RIR训练"。既有VBT实证研究的对照组普遍设定为传统\%1RM处方，其相较于ACSM 2026立场声明所推荐的规范化RIR训练是否仍能维持适应性优势，目前尚缺乏直接的随机对照证据。本研究通过在同一实验框架下直接比较两条已获科学认可的反馈路径，为抗阻训练自主调节理论在权威指南更新后的方法学再校准提供实证素材。第二，本研究采用了整合测量效度、干预效果与用户体验三个证据层次的评估思路，既有关于数字化训练工具的研究多聚焦于单一维度的测量精度验证，较少将技术工具置于随机对照干预中检验其相较于规范化传统方法的实际增量价值，也较少从用户体验层面追踪工具接受度的分化特征。本研究通过量化过程数据、量化结局数据与质性访谈数据的三角互证，为同类AI辅助训练工具的评估方案设计提供了一个具体的方法学案例。

\subsubsection{实践意义}

本研究的实践价值按由近及远的顺序体现在三个层次。在直接产出层面，本研究通过对自研移动端AI系统在真实训练场景中同步记录数据的分析，揭示了其相较于专业线性位置传感器的一致性水平与系统性偏差特征，为该系统的算法校准与后续版本改进提供了直接的工程参数。在指导实践层面，本研究客观呈现了客观速度反馈与规范化RIR反馈在实施成本、执行质量、体能适应与主观体验上的相对表现，为具有类似特征的训练者与体能训练从业者在权衡是否引入客观速度监控时提供了初步的实证参照。在方法学参数层面，本研究作为一项严格遵循预试验性研究规范开展的探索性工作，所报告的招募保留特征、系统可用性瓶颈、各项体能结局的效应量估计与训练执行过程的关键参数，可为后续开展大规模确证性随机对照试验的样本量计算、随访策略设计与干预方案完善提供直接的先验参数。
